% ! TEX root = ../m3-19-20.tex

\chapter{Scheme clasice de probabilitate}

Schemele clasice de probabilitate sînt modele matematice care ne permit calculul
probabilității de realizare a unui eveniment în cazul unor distribuții anume.

\section{Schema lui Poisson}

\indent\indent Schema lui Poisson se formulează astfel. Presupunem că avem $ n $
urne $ U_i, 1 \leq i \leq n $, care conțin bile albe și bile negre în proporții
cunoscute. Fie $ p_i $ probabilitatea de extragere a unei bile albe din urna $ U_i $,
respectiv $ q_i $, probabilitatea de extragere a unei bile negre din urna $ U_i $.

Extragem cîte o bilă din fiecare urnă. Probabilitatea de a obține $ k $ bile albe
este coeficientul lui $ X^k $ din polinomul:
\[
  \prod_{i = 1}^n (p_i X + q_i) = (p_1 X + q_1) (p_2 X + q_2) \cdots (p_n X + q_n).
\]

\textbf{Observație:} În cazul schemei lui Poisson, extragerea se face \emph{fără %
  revenire}. O bilă, după ce a fost extrasă, nu este repusă în urnă, astfel că, după
fiecare extragere, numărul de bile din urnă scade.

\vspace{0.3cm}

\textbf{Exemplu:} Avem 3 sisteme de siguranță, primul funcționează cu probabilitatea
de 80\%, al doilea, de 70\%, iar al treilea, în 90\% din cazuri. Să se determine
probabilitatea ca oricare două dintre sisteme să funcționeze simultan.

\emph{Soluție:} Ne aflăm în cazul schemei lui Poisson, \qq{bilele} fiind sistemele
de siguranță, iar \qq{extragerea} înseamnă funcționarea lor. Așadar, căutăm
coeficientul lui $ X^2 $ din polinomul:
\[
  \pi = (0,8 \cdot X + 0,2)(0,7 \cdot X + 0,3)(0,9 \cdot X + 0,1),
\]
coeficient care este $ p = 0,398 $.

%%%%%%%%%%%%%%%%%%%%%%%%%%%%%%%%%%%%%%%%%%%%%%%%%%%%%%%%%%%%%%%%%%%%%%

\section{Schema lui Bernoulli (binomială)}

\indent\indent Situația este similară cu cea din cazul Poisson, doar că acum
extragerile se fac \emph{cu repunere}. Adică, după ce fiecare bilă este extrasă
și se înregistrează culoarea sa, este repusă în urnă, pentru a putea fi extrasă din
nou, eventual.

Așadar, schema lui Bernoulli se formulează astfel. Avem o urnă cu $ a $ bile albe
și $ b $ bile negre. Extragem cu repunere $ n $ bile. Probabilitatea să extragem
$ k $ bile albe, cu $ 0 \leq k \leq n $ este dată de:
\[
  p(n, k) = C_n^k p^k q^{n-k}, \quad \text{unde } p = \frac{a}{a + b}, %
  q = \frac{b}{a + b}.
\]

De remarcat că în acest caz, $ p $ corespunde evenimentului favorabil, iar $ q = 1 - p $.

\vspace{0.3cm}

\textbf{Exemplu:} Se aruncă două zaruri de 10 ori. Care este probabilitatea să se obțină
de exact 4 ori suma 7?

\emph{Soluție:} Ne aflăm în cazul schemei lui Bernoulli, unde \qq{bilele} sînt sumele
de pe zaruri, iar \qq{extragerea} este aruncarea zarului. Evident, modelul potrivit
este acela al extragerii cu revenire, deoarece orice sumă poate fi obținută la fiecare
aruncare.

În acest caz, avem $ p = \dfrac{1}{6} $, deoarece suma 7 poate fi obținută în 6 moduri
din punctele de pe două zaruri, iar cazurile posibile sînt 36. Rezultă $ q = \dfrac{5}{6} $.

În plus, avem $ n = 10, k = 4 $, deci:
\[
  p(10, 4) = C_{10}^4 \frac{1}{6^4} \cdot \frac{5^6}{6^6}.
\]

%%%%%%%%%%%%%%%%%%%%%%%%%%%%%%%%%%%%%%%%%%%%%%%%%%%%%%%%%%%%%%%%%%%%%%

\section{Schema multinomială}

\indent\indent Schema lui Bernoulli se mai numește \emph{schema binomială},
deoarece $ p(n, k) $ este, de fapt, coeficientul lui $ X^k $ din dezvoltarea
binomului $ (pX + q)^n $.

O generalizare a acestei situații este \emph{schema multinomială}, în care
presupunem că avem o urnă cu bile de $ s \in \NN $ culori și se extrag cu
repunere $ n $ bile. Probabilitatea de a extrage $ k_i $ bile de culoare
$ i $, cu $ 1 \leq i \leq s $ este dată de formula:
\[
  p(n; k_1, \dots, k_s) = \frac{n!}{k_1! k_2! \cdots k_s!} \cdot %
  p_1^{k_1} \cdot p_2^{k_2} \cdots p_s^{k_s},
\]
unde $ n = k_1 + \dots + k_s $ și $ p_1 + \dots + p_s = 1 $, iar
$ p_i $ este probabilitatea de a extrage o bilă de culoare $ i $.

\vspace{0.3cm}

\textbf{Exemplu:} Se aruncă cu un zar de 10 ori. Care este probabilitatea
ca exact de 2 ori să apară fața cu 1 punct și de 3 ori să apară fața cu 2 puncte?

\emph{Soluție:} Evident, sîntem în cazul schemei multinomiale. \qq{Culorile}
sînt punctele de pe zaruri și vrem să \qq{extragem 2 bile de culoarea 1 și 3
  bile de culoarea 2}.

Atunci $ n = 10 $, deoarece facem 10 \qq{extrageri}, $ k_1 = 2 $ și $ k_2 = 3 $.
În plus, avem $ p_1 = p_2 = \dfrac{1}{6} $, deoarece orice număr de pe zar are
aceeași probabilitate de a apărea. Rezultă, de asemenea, că avem nevoie și
de $ k_3 = n - k_1 - k_2 = 5 $ și $ p_3 = 1 - p_1 - p_2 = \frac{2}{3} $.

Probabilitatea cerută este:
\[
  p(10; 2, 3, 5) = \frac{10!}{2! \cdot 3! \cdot 5!} \cdot %
  \frac{1}{6^2} \cdot \frac{1}{6^3} \cdot \frac{2^5}{3^5}.
\]

%%%%%%%%%%%%%%%%%%%%%%%%%%%%%%%%%%%%%%%%%%%%%%%%%%%%%%%%%%%%%%%%%%%%%%

\section{Schema geometrică}

\indent\indent Dintr-o urnă cu $ a $ bile albe și $ b $ bile negre, extragem
fără repunere $ n $ bile, cu $ n \leq a + b $. Probabilitatea de a obține
$ k \leq a $ bile albe, \emph{fără repunere}, este dată de:
\[
  p_{a, b}^{k, n - k} = \frac{C_a^k \cdot C_b^{n-k}}{C_{a+b}^n}.
\]

Mai general, dacă avem $ a_i $ bile de culoarea $ i, 1 \leq i \leq s $ și
extragem $ n $ bile fără repunere, atunci probabilitatea de a extrage $ k_i $,
cu $ k_1 + \dots + k_s = n $ bile de culoarea $ i $ este dată de formula:
\[
  p_{a_1, \dots, a_s}^{k_1, \dots, k_s} = %
  \frac{C_{a_1}^{k_1} \cdots C_{a_s}^{k_s}}{C_{k_1 + \dots + k_s}^n}.
\]

\vspace{0.3cm}

\textbf{Exemplu:} Avem un lot de 100 articole, dintre care 80 sînt corespunzătoare,
15 sînt cu defecțiuni remediabile și 5 rebuturi. Alegem 6 articole. Care este
probabilitatea ca dintre acestea, 3 să fie bune, 2 cu defecțiuni remediabile și
un rebut?

\emph{Soluție:} Presupunem că extragerile se fac cu repunere. Atunci sîntem în cazul
schemei multinomiale, deci:
\[
  p(6; 3, 2, 1) = \frac{6!}{3! \cdot 2! \cdot 1!} \Big(\frac{80}{100}\Big)^3 \cdot %
  \Big( \frac{15}{100} \Big)^2 \cdot \Big( \frac{5}{100} \Big)^1.
\]

Dacă extragerile se fac fără repunere, ne aflăm în cazul schemei geometrice și avem:
\[
  p_{80, 15, 5}^{3, 2, 1} = \frac{C_{80}^3 \cdot C_{15}^2 \cdot C_5^1}{C_{100}^6}.
\]


%%%%%%%%%%%%%%%%%%%%%%%%%%%%%%%%%%%%%%%%%%%%%%%%%%%%%%%%%%%%%%%%%%%%%%

\section{Exerciții}

\indent\indent 1. Un lot de 50 de procesoare conține 3 defecte. Alegem la întîmplare
10, simultan, și le verificăm. Care este probabilitatea ca 2 să fie defecte?

\emph{Indicație:} Schema geometrică (fără repunere).

\vspace{0.3cm}

2. (a) Un semnal este transmis pe 3 canale diferite, iar probabilitățile de
recepționare corectă sînt 0,9; 0,8 și 0,7. Care este probabilitatea să se
recepționeze corect un semnal?

(b) Dar dacă semnalul este transmis pe un canal ales la întîmplare și presupunem
că orice canal poate fi ales cu aceeași probabilitate?

\emph{Indicație:} (a) Schema lui Poisson.

(b) Formula probabilității totale: $ \dfrac{1}{3} \cdot (0,9 + 0,8 + 0,7) $.

\vspace{0.3cm}

3. O urnă conține 2 bile albe și 3 bile negre. Alegem la întîmplare 2 bile,
fără repunere. Care este probabilitatea ca acestea să fie ambele negre?

Dar dacă extragerea se face cu repunere?

\vspace{0.3cm}

4. Se testează 5 dispozitive care funcționează în condiții identice, independent
și cu randament de 0,9 fiecare. Care este probabilitatea ca exact două să
funcționeze?

\emph{Indicație:} Schema Bernoulli.

\vspace{0.3cm}

5. Se consideră 3 urne: $ U_1 $, care conține 5 bile albe și 5 negre, $ U_2 $,
care conține 4 bile albe și 6 negre și $ U_3 $, care conține 4 bile albe
și 5 negre. Din fiecare urnă se extrag cu repunere cîte 5 bile.

Care este probabilitatea ca din 2 urne să obținem cîte 2 bile albe și 3 negre,
iar din a treia urnă să obținem o altă combinație?

\emph{Indicație:} Schema Poisson sau Bernoulli, după cazuri (cu/fără repunere).

\vspace{0.3cm}

6. Se consideră urnele din problema precedentă. Din fiecare urnă se extrage
cîte o bilă. Dacă se repetă experiența de 5 ori, care este probabilitatea ca
de 3 ori să se obțină o bilă albă și 2 negre?

\emph{Indicație:} Schema Poisson.

\vspace{0.3cm}

7. Fie 8 canale de transmisie a informației care funcționează independent.
Presupunem că un canal este activ cu probabilitatea 1/3.

Să se calculeze probabilitatea ca la un moment dat să fie mai mult de 6 canale
active.

\emph{Indicație:} Schema Bernoulli.

\vspace{0.3cm}

8. Șase vînători au zărit o vulpe și au tras simultan. Presupunem că de la
distanța de la care au tras, fiecare vînător nimerește în mod obișnuit vulpea
cu o probabilitate de $ \dfrac{1}{3} $. Aflați probabilitatea ca vulpea să
fie nimerită.

\vspace{0.3cm}

9. Se aruncă o monedă de 8 ori. Aflați probabilitatea ca stema să apară
de 6 ori. Aflați probabilitatea ca stema să apară de \emph{cel puțin} 6 ori.

\vspace{0.3cm}

10. Un muncitor produce piese astfel încît 99\% sînt bune, 7\% au defecte
remediabile și 3\% sînt rebuturi.

Se aleg aleatoriu 3 piese produse de muncitor. Care este probabilitatea
ca dintre aceste 3 piese, cel puțin una să fie bună și cel puțin una să fie
rebut?


%%% Local Variables:
%%% mode: latex
%%% TeX-master: "../m3-19-20"
%%% End:
